% ============================================================
%  sections/chapter_02.tex
%  Section 2: Mathematical Background
% ============================================================

\section{Mathematical Background}
\label{sec:background}

This section establishes the notation and mathematical foundations used
throughout the report.
We introduce tensor algebra, the two decomposition families under study
(Tucker and CP), the multilinear singular value decomposition (HOSVD),
and the reconstruction losses that define each optimization problem.

% ------------------------------------------------------------
\subsection{Tensor Notation}
\label{subsec:notation}
% ------------------------------------------------------------

A \emph{tensor} $\tensor{X} \in \R^{I_1 \times I_2 \times \cdots \times I_N}$
is an $N$-way array.
Its entry indexed by $(i_1, i_2, \ldots, i_N)$ is written $x_{i_1 i_2 \cdots i_N}$.
We adopt the following notational convention throughout:
calligraphic letters ($\tensor{X}$) for tensors,
bold uppercase ($\mat{U}$) for matrices,
bold lowercase ($\vect{u}$) for vectors,
and italic lowercase ($\lambda, r$) for scalars.

\paragraph{Frobenius norm.}
The Frobenius norm of a tensor generalizes the matrix Frobenius norm
entry-wise:
\begin{equation}
  \frobn{\tensor{X}}
  \;=\; \sqrt{\,\sum_{i_1=1}^{I_1} \cdots \sum_{i_N=1}^{I_N}
               x_{i_1 \cdots i_N}^2\,}
  \;=\; \sqrt{\langle \tensor{X},\, \tensor{X} \rangle}.
  \label{eq:frobenius}
\end{equation}
It satisfies \emph{unitary invariance}: for any orthogonal matrices
$\mat{Q}_n$,
\begin{equation}
  \frobn{\tensor{X} \times_1 \mat{Q}_1 \times_2 \mat{Q}_2 \cdots
         \times_N \mat{Q}_N}
  \;=\; \frobn{\tensor{X}}.
  \label{eq:unitary_inv}
\end{equation}
This property is central to the correctness of HOSVD and HOOI
(Section~\ref{subsec:hosvd}).

\paragraph{Mode-$n$ unfolding (matricization).}
The \emph{mode-$n$ unfolding}
$\unfold{X}{n} \in \R^{I_n \times \prod_{k \neq n} I_k}$
rearranges the entries of $\tensor{X}$ such that the mode-$n$ index
maps to rows and all remaining indices (in lexicographic order) map to
columns~\cite{kolda2009tensor}.
Concretely,
\begin{equation}
  \bigl[\unfold{X}{n}\bigr]_{i_n,\, j}
  \;=\; x_{i_1 i_2 \cdots i_N},
  \qquad
  j = 1 + \sum_{\substack{k=1\\k\neq n}}^{N}
          (i_k - 1)\prod_{\substack{m=1\\m\neq n}}^{k-1} I_m.
  \label{eq:unfolding}
\end{equation}

\paragraph{Mode-$n$ product.}
The \emph{mode-$n$ product} of
$\tensor{X} \in \R^{I_1 \times \cdots \times I_N}$
with $\mat{U} \in \R^{J \times I_n}$ is the tensor
$\tensor{Y} = \tensor{X} \times_n \mat{U}
\in \R^{I_1 \times \cdots \times I_{n-1} \times J \times I_{n+1}
\times \cdots \times I_N}$ defined entry-wise as
\begin{equation}
  y_{i_1 \cdots i_{n-1}\, j\, i_{n+1} \cdots i_N}
  \;=\; \sum_{i_n=1}^{I_n} x_{i_1 \cdots i_N}\, u_{j i_n}.
  \label{eq:mode_product}
\end{equation}
In matricized form this becomes
\begin{equation}
  \bigl(\tensor{X} \times_n \mat{U}\bigr)_{(n)}
  \;=\; \mat{U}\, \unfold{X}{n},
  \label{eq:mode_product_mat}
\end{equation}
which reduces the mode-$n$ product to ordinary matrix multiplication and
underpins the ALS update rules derived in
Section~\ref{subsec:als}.

\paragraph{Commutativity of mode products.}
For $m \neq n$,
\begin{equation}
  \tensor{X} \times_m \mat{A} \times_n \mat{B}
  \;=\; \tensor{X} \times_n \mat{B} \times_m \mat{A}.
  \label{eq:mode_commute}
\end{equation}
This commutativity is used repeatedly in the derivation of ALS subproblems.

% ------------------------------------------------------------
\subsection{Tucker Decomposition}
\label{subsec:tucker}
% ------------------------------------------------------------

\begin{definition}[Tucker decomposition]
  A rank-$(r_1, \ldots, r_N)$ Tucker approximation of
  $\tensor{X} \in \R^{I_1 \times \cdots \times I_N}$ is
  \begin{equation}
    \hat{\tensor{X}}
    \;=\; \tensor{G}
          \times_1 \mat{U}^{(1)}
          \times_2 \mat{U}^{(2)}
          \cdots
          \times_N \mat{U}^{(N)},
    \label{eq:tucker}
  \end{equation}
  where $\tensor{G} \in \R^{r_1 \times \cdots \times r_N}$ is the
  \emph{core tensor} and each factor matrix
  $\mat{U}^{(n)} \in \R^{I_n \times r_n}$ has orthonormal columns,
  i.e., ${\mat{U}^{(n)}}^\top \mat{U}^{(n)} = \mat{I}_{r_n}$.
\end{definition}

Mode-$n$ unfolding of $\hat{\tensor{X}}$ yields
\begin{equation}
  \hat{\mat{X}}_{(n)}
  \;=\; \mat{U}^{(n)}\, \mat{G}_{(n)}
        \Bigl(\mat{U}^{(N)} \kron \cdots \kron \mat{U}^{(n+1)}
              \kron \mat{U}^{(n-1)} \kron \cdots \kron \mat{U}^{(1)}
        \Bigr)^{\!\top},
  \label{eq:tucker_matricized}
\end{equation}
where $\kron$ denotes the Kronecker product.
Equation~\eqref{eq:tucker_matricized} exposes the least-squares structure
exploited by ALS.

\paragraph{Parameter count and compression ratio.}
The Tucker format requires
\begin{equation}
  P_{\mathrm{Tucker}}
  \;=\; \prod_{n=1}^{N} r_n
        \;+\; \sum_{n=1}^{N} I_n r_n
  \label{eq:tucker_params}
\end{equation}
parameters.
The compression ratio relative to the dense tensor is
\begin{equation}
  \rho
  \;=\; \frac{\displaystyle\prod_{n=1}^{N} I_n}
             {P_{\mathrm{Tucker}}}.
  \label{eq:tucker_ratio}
\end{equation}
For a three-channel image of size $H \times W \times 3$ with Tucker rank
$(r, r, 3)$, this gives
$\rho = HW \cdot 3\,/\,(3r^2 + Hr + Wr + 3r)$,
which grows roughly as $HW/r^2$ for large images.

\paragraph{Relationship to truncated SVD.}
For $N=2$, Tucker reduces to the truncated SVD:
$\hat{\mat{X}} = \mat{U} \mat{\Sigma}_r \mat{V}^\top$,
with $\tensor{G} = \mat{\Sigma}_r$ and
$\mat{U}^{(1)} = \mat{U}$, $\mat{U}^{(2)} = \mat{V}$.
Tucker thus provides the natural multilinear generalization of low-rank
matrix approximation.

% ------------------------------------------------------------
\subsection{CP Decomposition}
\label{subsec:cp}
% ------------------------------------------------------------

\begin{definition}[CP decomposition]
  A rank-$R$ CANDECOMP/PARAFAC (CP) approximation of $\tensor{X}$ is
  \begin{equation}
    \hat{\tensor{X}}
    \;=\; \sum_{r=1}^{R} \lambda_r\,
          \vect{a}_r^{(1)} \circ \vect{a}_r^{(2)} \circ \cdots \circ
          \vect{a}_r^{(N)},
    \label{eq:cp}
  \end{equation}
  where $\circ$ denotes the outer product,
  $\|\vect{a}_r^{(n)}\|_2 = 1$, and $\lambda_r \geq 0$.
\end{definition}

CP can be viewed as Tucker with a \emph{superdiagonal} core of size
$R \times \cdots \times R$, i.e., the case $r_1 = \cdots = r_N = R$
with $g_{i_1 \cdots i_N} = \lambda_{i_1}$ only when
$i_1 = i_2 = \cdots = i_N$ and zero otherwise.
The CP format stores
$R \bigl(\sum_{n=1}^{N} I_n + 1\bigr)$
parameters, fewer than Tucker for the same $R$, but determining the exact
CP rank is NP-hard in general~\cite{kolda2009tensor}.

\begin{remark}[CP degeneracy]
  Unlike Tucker, the best rank-$R$ CP approximation of a tensor of order
  $N \geq 3$ may not exist: the infimum of the objective is not always
  attained~\cite{kolda2009tensor}.
  Tucker decomposition avoids this pathology because orthonormal factor
  matrices form a compact set.
\end{remark}

% ------------------------------------------------------------
\subsection{Higher-Order SVD (HOSVD)}
\label{subsec:hosvd}
% ------------------------------------------------------------

\begin{definition}[HOSVD~\cite{de2000multilinear}]
  The \emph{Higher-Order SVD} of $\tensor{X}$ computes each factor matrix
  independently:
  \begin{equation}
    \mat{U}^{(n)}
    \;\coloneqq\;
    \text{leading } r_n \text{ left singular vectors of } \unfold{X}{n},
    \quad n = 1, \ldots, N,
    \label{eq:hosvd_factors}
  \end{equation}
  and sets the core tensor to
  \begin{equation}
    \tensor{G}
    \;=\; \tensor{X}
          \times_1 {\mat{U}^{(1)}}^\top
          \times_2 {\mat{U}^{(2)}}^\top
          \cdots
          \times_N {\mat{U}^{(N)}}^\top.
    \label{eq:hosvd_core}
  \end{equation}
\end{definition}

HOSVD is computed in closed form and provides a deterministic initialization.
However, it is \emph{not} the globally optimal Tucker approximation because
the mode-$n$ factor matrices are optimized independently rather than jointly.

\begin{proposition}[HOSVD error bound~\cite{de2000multilinear}]
  Let $\hat{\tensor{X}}_{\mathrm{HOSVD}}$ denote the rank-$(r_1,\ldots,r_N)$
  HOSVD approximation of $\tensor{X}$.
  Then
  \begin{equation}
    \frobn{\tensor{X} - \hat{\tensor{X}}_{\mathrm{HOSVD}}}^2
    \;\leq\;
    \sum_{n=1}^{N} \sum_{i=r_n+1}^{I_n} \sigma_i^{(n)2},
    \label{eq:hosvd_bound}
  \end{equation}
  where $\sigma_i^{(n)}$ denotes the $i$-th singular value of $\unfold{X}{n}$.
\end{proposition}

Equation~\eqref{eq:hosvd_bound} is used in Section~\ref{sec:experiments}
to benchmark the rank sweep experiment:
the left-hand side is measured empirically and compared with the right-hand
side computed from the singular value spectra.

% ------------------------------------------------------------
\subsection{Reconstruction Losses}
\label{subsec:losses}
% ------------------------------------------------------------

The general approximation problem is
\begin{equation}
  \min_{\tensor{G},\, \{\mat{U}^{(n)}\}}\;
  \mathcal{L}\!\left(\tensor{X},\, \hat{\tensor{X}}\right),
  \label{eq:general_opt}
\end{equation}
where the choice of $\mathcal{L}$ determines both the statistical
interpretation of the solution and the algorithmic complexity of
the optimization.
We study three loss functions.

\paragraph{Squared Frobenius loss.}
\begin{equation}
  \mathcal{L}_F(\tensor{X}, \hat{\tensor{X}})
  \;=\; \frobn{\tensor{X} - \hat{\tensor{X}}}^2
  \;=\; \sum_{i_1, \ldots, i_N}
        \bigl(x_{i_1\cdots i_N} - \hat{x}_{i_1\cdots i_N}\bigr)^2.
  \label{eq:loss_frobenius}
\end{equation}
This loss is smooth and convex in each factor separately,
enabling closed-form ALS updates.
When the noise is i.i.d.\ Gaussian, the Frobenius minimizer is the
maximum likelihood estimator.
Its sensitivity to large residuals, however, makes it non-robust to outliers.

\paragraph{$\ell_1$ loss.}
\begin{equation}
  \mathcal{L}_1(\tensor{X}, \hat{\tensor{X}})
  \;=\; \onenorm{\tensor{X} - \hat{\tensor{X}}}
  \;=\; \sum_{i_1, \ldots, i_N}
        \bigl|x_{i_1\cdots i_N} - \hat{x}_{i_1\cdots i_N}\bigr|.
  \label{eq:loss_l1}
\end{equation}
Minimizing $\mathcal{L}_1$ corresponds to maximum likelihood estimation
under a Laplace noise model and yields a positive breakdown point,
making it robust to sparse outliers.
The loss is convex but \emph{non-smooth} at zero; its subdifferential is
$\partial|\cdot| = \sign(\cdot)$, which precludes direct use of
standard least-squares solvers.

\paragraph{Huber loss.}
The Huber loss~\cite{huber1964robust} interpolates smoothly between
$\ell_2$ (near zero) and $\ell_1$ (in the tails):
\begin{equation}
  \rho_\delta(r)
  \;=\; \begin{cases}
          \dfrac{1}{2} r^2
            & \text{if } |r| \leq \delta, \\[8pt]
          \delta \!\left(|r| - \dfrac{\delta}{2}\right)
            & \text{if } |r| > \delta,
        \end{cases}
  \label{eq:huber}
\end{equation}
with threshold $\delta > 0$.
The aggregate loss is
$\mathcal{L}_H = \sum_{i_1,\ldots,i_N} \rho_\delta(x_{i_1\cdots i_N}
- \hat{x}_{i_1\cdots i_N})$.
Being everywhere differentiable, $\mathcal{L}_H$ admits gradient-based
optimization without subgradient complications.

\paragraph{Summary.}
Table~\ref{tab:losses} summarizes the three losses with respect to
statistical interpretation, smoothness, and optimization strategy.

\begin{table}[h]
  \centering
  \caption{Comparison of reconstruction losses.}
  \label{tab:losses}
  \begin{tabular}{lccc}
    \toprule
    Loss & Noise model & Smoothness & Optimization \\
    \midrule
    $\mathcal{L}_F$ (Frobenius)
      & Gaussian      & $C^\infty$       & ALS (closed-form) \\
    $\mathcal{L}_1$ ($\ell_1$)
      & Laplace       & non-smooth at 0  & IRLS / subgradient \\
    $\mathcal{L}_H$ (Huber)
      & Gaussian + outliers & $C^1$      & Gradient descent \\
    \bottomrule
  \end{tabular}
\end{table}

% % ============================================================
% %  sections/chapter_02_ko.tex
% %  2장: 수학적 배경
% % ============================================================

% \section{수학적 배경}
% \label{sec:background}

% 이 장에서는 보고서 전반에 걸쳐 사용되는 표기법과 수학적 기초를 정립한다.
% 텐서 대수, 두 가지 분해 방법(Tucker 및 CP),
% 고차 특이값 분해(HOSVD), 그리고 각 최적화 문제를 정의하는
% 재구성 손실 함수를 순서대로 소개한다.

% % ------------------------------------------------------------
% \subsection{텐서 표기법}
% \label{subsec:notation}
% % ------------------------------------------------------------

% 텐서(tensor) $\tensor{X} \in \R^{I_1 \times I_2 \times \cdots \times I_N}$는
% $N$-방향(way) 배열이며, 인덱스 $(i_1, i_2, \ldots, i_N)$에 해당하는
% 원소를 $x_{i_1 i_2 \cdots i_N}$으로 표기한다.
% 본 보고서에서는 다음의 표기 규약을 일관되게 사용한다:
% 캘리그래피 문자($\tensor{X}$)는 텐서, 굵은 대문자($\mat{U}$)는 행렬,
% 굵은 소문자($\vect{u}$)는 벡터, 이탤릭체 소문자($\lambda, r$)는 스칼라를 나타낸다.

% \paragraph{프로베니우스 노름(Frobenius norm).}
% 텐서의 프로베니우스 노름은 행렬의 경우를 원소별로 일반화한 것이다:
% \begin{equation}
%   \frobn{\tensor{X}}
%   \;=\; \sqrt{\,\sum_{i_1=1}^{I_1} \cdots \sum_{i_N=1}^{I_N}
%                x_{i_1 \cdots i_N}^2\,}
%   \;=\; \sqrt{\langle \tensor{X},\, \tensor{X} \rangle}.
%   \label{eq:frobenius}
% \end{equation}
% 프로베니우스 노름은 \emph{유니터리 불변성(unitary invariance)}을 만족한다.
% 임의의 직교 행렬 $\mat{Q}_n$에 대해,
% \begin{equation}
%   \frobn{\tensor{X} \times_1 \mat{Q}_1 \times_2 \mat{Q}_2 \cdots
%          \times_N \mat{Q}_N}
%   \;=\; \frobn{\tensor{X}}.
%   \label{eq:unitary_inv}
% \end{equation}
% 이 성질은 \ref{subsec:hosvd}절에서 설명하는 HOSVD와 HOOI의
% 정확성을 보장하는 핵심 근거가 된다.

% \paragraph{모드-$n$ 전개(mode-$n$ unfolding, matricization).}
% \emph{모드-$n$ 전개} $\unfold{X}{n} \in \R^{I_n \times \prod_{k \neq n} I_k}$는
% $\tensor{X}$의 원소를 모드-$n$ 인덱스가 행(row)에,
% 나머지 인덱스가 사전 순서(lexicographic order)에 따라 열(column)에
% 배치되도록 재배열한 것이다~\cite{kolda2009tensor}.
% 구체적으로,
% \begin{equation}
%   \bigl[\unfold{X}{n}\bigr]_{i_n,\, j}
%   \;=\; x_{i_1 i_2 \cdots i_N},
%   \qquad
%   j = 1 + \sum_{\substack{k=1\\k\neq n}}^{N}
%           (i_k - 1)\prod_{\substack{m=1\\m\neq n}}^{k-1} I_m.
%   \label{eq:unfolding}
% \end{equation}

% \paragraph{모드-$n$ 곱(mode-$n$ product).}
% $\tensor{X} \in \R^{I_1 \times \cdots \times I_N}$와
% $\mat{U} \in \R^{J \times I_n}$의 \emph{모드-$n$ 곱}은
% \begin{equation}
%   y_{i_1 \cdots i_{n-1}\, j\, i_{n+1} \cdots i_N}
%   \;=\; \sum_{i_n=1}^{I_n} x_{i_1 \cdots i_N}\, u_{j i_n}
%   \label{eq:mode_product}
% \end{equation}
% 으로 정의되는 텐서
% $\tensor{Y} = \tensor{X} \times_n \mat{U}
% \in \R^{I_1 \times \cdots \times I_{n-1} \times J
% \times I_{n+1} \times \cdots \times I_N}$이다.
% 행렬 형태로 나타내면,
% \begin{equation}
%   \bigl(\tensor{X} \times_n \mat{U}\bigr)_{(n)}
%   \;=\; \mat{U}\, \unfold{X}{n},
%   \label{eq:mode_product_mat}
% \end{equation}
% 즉 모드-$n$ 곱은 일반적인 행렬 곱셈으로 환원된다.
% 이 관계식은 \ref{subsec:als}절의 ALS 갱신 규칙을 유도하는 기초가 된다.

% \paragraph{모드 곱의 교환 법칙.}
% $m \neq n$인 경우,
% \begin{equation}
%   \tensor{X} \times_m \mat{A} \times_n \mat{B}
%   \;=\; \tensor{X} \times_n \mat{B} \times_m \mat{A}.
%   \label{eq:mode_commute}
% \end{equation}
% 이 교환 법칙은 ALS 부분 문제(subproblem) 유도에 반복적으로 사용된다.

% % ------------------------------------------------------------
% \subsection{Tucker 분해}
% \label{subsec:tucker}
% % ------------------------------------------------------------

% \begin{definition}[Tucker 분해]
%   $\tensor{X} \in \R^{I_1 \times \cdots \times I_N}$의
%   랭크-$(r_1, \ldots, r_N)$ Tucker 근사는
%   \begin{equation}
%     \hat{\tensor{X}}
%     \;=\; \tensor{G}
%           \times_1 \mat{U}^{(1)}
%           \times_2 \mat{U}^{(2)}
%           \cdots
%           \times_N \mat{U}^{(N)}
%     \label{eq:tucker}
%   \end{equation}
%   이며, 여기서 $\tensor{G} \in \R^{r_1 \times \cdots \times r_N}$는
%   \emph{코어 텐서(core tensor)},
%   각 $\mat{U}^{(n)} \in \R^{I_n \times r_n}$는
%   열이 정규직교(orthonormal)인 \emph{인수 행렬(factor matrix)}이다.
%   즉 ${\mat{U}^{(n)}}^\top \mat{U}^{(n)} = \mat{I}_{r_n}$이다.
% \end{definition}

% $\hat{\tensor{X}}$의 모드-$n$ 전개는 다음과 같다:
% \begin{equation}
%   \hat{\mat{X}}_{(n)}
%   \;=\; \mat{U}^{(n)}\, \mat{G}_{(n)}
%         \Bigl(\mat{U}^{(N)} \kron \cdots \kron \mat{U}^{(n+1)}
%               \kron \mat{U}^{(n-1)} \kron \cdots \kron \mat{U}^{(1)}
%         \Bigr)^{\!\top},
%   \label{eq:tucker_matricized}
% \end{equation}
% 여기서 $\kron$은 크로네커 곱(Kronecker product)이다.
% 식~\eqref{eq:tucker_matricized}는 ALS가 활용하는 최소제곱 구조를 명시적으로 드러낸다.

% \paragraph{파라미터 수와 압축률.}
% Tucker 형식이 저장하는 파라미터 수는
% \begin{equation}
%   P_{\mathrm{Tucker}}
%   \;=\; \prod_{n=1}^{N} r_n
%         \;+\; \sum_{n=1}^{N} I_n r_n
%   \label{eq:tucker_params}
% \end{equation}
% 이며, 원본 텐서 대비 압축률은
% \begin{equation}
%   \rho
%   \;=\; \frac{\displaystyle\prod_{n=1}^{N} I_n}
%              {P_{\mathrm{Tucker}}}
%   \label{eq:tucker_ratio}
% \end{equation}
% 이다.
% 크기 $H \times W \times 3$인 컬러 이미지에 Tucker 랭크 $(r, r, 3)$을
% 적용하면 $\rho = 3HW\,/\,(3r^2 + Hr + Wr + 3r)$이며,
% 대형 이미지에서 $r$이 작을수록 $\rho \approx HW/r^2$에 가까워진다.

% \paragraph{절단 SVD와의 관계.}
% $N=2$이면 Tucker는 절단 SVD로 환원된다:
% $\hat{\mat{X}} = \mat{U} \mat{\Sigma}_r \mat{V}^\top$,
% $\tensor{G} = \mat{\Sigma}_r$,
% $\mat{U}^{(1)} = \mat{U}$, $\mat{U}^{(2)} = \mat{V}$.
% Tucker 분해는 저랭크 행렬 근사의 자연스러운 다중선형(multilinear) 일반화이다.

% % ------------------------------------------------------------
% \subsection{CP 분해}
% \label{subsec:cp}
% % ------------------------------------------------------------

% \begin{definition}[CP 분해]
%   $\tensor{X}$의 랭크-$R$ CANDECOMP/PARAFAC (CP) 근사는
%   \begin{equation}
%     \hat{\tensor{X}}
%     \;=\; \sum_{r=1}^{R} \lambda_r\,
%           \vect{a}_r^{(1)} \circ \vect{a}_r^{(2)} \circ \cdots \circ
%           \vect{a}_r^{(N)}
%     \label{eq:cp}
%   \end{equation}
%   이며, 여기서 $\circ$는 외적(outer product),
%   $\|\vect{a}_r^{(n)}\|_2 = 1$, $\lambda_r \geq 0$이다.
% \end{definition}

% CP는 코어 텐서가 $R \times \cdots \times R$ 크기의
% \emph{초대각(superdiagonal)} 텐서로 제한된 Tucker 분해로 볼 수 있다.
% CP 형식은 $R \bigl(\sum_{n=1}^{N} I_n + 1\bigr)$개의 파라미터를 저장하므로
% 동일한 $R$에서 Tucker보다 더 간결하지만,
% 정확한 CP 랭크를 결정하는 문제는 일반적으로 NP-난해(NP-hard)이다~\cite{kolda2009tensor}.

% \begin{remark}[CP 퇴화(CP degeneracy)]
%   Tucker와 달리, 차수 $N \geq 3$인 텐서의 최적 랭크-$R$ CP 근사는
%   존재하지 않을 수 있다 --- 목적 함수의 하한이 항상 달성되지는 않는다~\cite{kolda2009tensor}.
%   Tucker 분해는 인수 행렬을 정규직교로 제약함으로써 이 문제를 회피한다.
% \end{remark}

% % ------------------------------------------------------------
% \subsection{고차 특이값 분해 (HOSVD)}
% \label{subsec:hosvd}
% % ------------------------------------------------------------

% \begin{definition}[HOSVD~\cite{de2000multilinear}]
%   $\tensor{X}$의 \emph{고차 특이값 분해(Higher-Order SVD, HOSVD)}는
%   각 인수 행렬을 독립적으로 계산한다:
%   \begin{equation}
%     \mat{U}^{(n)}
%     \;\coloneqq\;
%     \unfold{X}{n}\text{의 주요 } r_n\text{개 좌특이벡터},
%     \quad n = 1, \ldots, N,
%     \label{eq:hosvd_factors}
%   \end{equation}
%   코어 텐서는
%   \begin{equation}
%     \tensor{G}
%     \;=\; \tensor{X}
%           \times_1 {\mat{U}^{(1)}}^\top
%           \times_2 {\mat{U}^{(2)}}^\top
%           \cdots
%           \times_N {\mat{U}^{(N)}}^\top
%     \label{eq:hosvd_core}
%   \end{equation}
%   으로 설정한다.
% \end{definition}

% HOSVD는 닫힌 형태(closed form)로 계산되며 결정론적 초기화를 제공한다.
% 그러나 각 모드의 인수 행렬을 독립적으로 최적화하므로
% 전역 최적 Tucker 근사가 아니다.

% \begin{proposition}[HOSVD 오차 상한~\cite{de2000multilinear}]
%   랭크-$(r_1, \ldots, r_N)$ HOSVD 근사 $\hat{\tensor{X}}_{\mathrm{HOSVD}}$에 대해,
%   \begin{equation}
%     \frobn{\tensor{X} - \hat{\tensor{X}}_{\mathrm{HOSVD}}}^2
%     \;\leq\;
%     \sum_{n=1}^{N} \sum_{i=r_n+1}^{I_n} \sigma_i^{(n)2},
%     \label{eq:hosvd_bound}
%   \end{equation}
%   여기서 $\sigma_i^{(n)}$는 $\unfold{X}{n}$의 $i$번째 특이값이다.
% \end{proposition}

% 식~\eqref{eq:hosvd_bound}은 \ref{sec:experiments}절의 랭크 스윕(rank sweep)
% 실험에서 이론적 기준값으로 사용된다.
% 실험에서 측정한 좌변과, 특이값 스펙트럼으로 계산한 우변을 비교하여
% 이론 예측의 정확도를 검증한다.

% % ------------------------------------------------------------
% \subsection{재구성 손실 함수}
% \label{subsec:losses}
% % ------------------------------------------------------------

% 일반적인 근사 문제는
% \begin{equation}
%   \min_{\tensor{G},\, \{\mat{U}^{(n)}\}}\;
%   \mathcal{L}\!\left(\tensor{X},\, \hat{\tensor{X}}\right)
%   \label{eq:general_opt}
% \end{equation}
% 로 표현된다.
% 손실 함수 $\mathcal{L}$의 선택은 해의 통계적 해석과 최적화 알고리즘의
% 복잡도를 동시에 결정한다.
% 본 연구에서는 세 가지 손실 함수를 비교한다.

% \paragraph{프로베니우스 제곱 손실.}
% \begin{equation}
%   \mathcal{L}_F(\tensor{X}, \hat{\tensor{X}})
%   \;=\; \frobn{\tensor{X} - \hat{\tensor{X}}}^2
%   \;=\; \sum_{i_1, \ldots, i_N}
%         \bigl(x_{i_1\cdots i_N} - \hat{x}_{i_1\cdots i_N}\bigr)^2.
%   \label{eq:loss_frobenius}
% \end{equation}
% 이 손실은 각 인수에 대해 매끄럽고(smooth) 볼록(convex)하여
% ALS의 닫힌 형태 갱신을 가능하게 한다.
% 잡음이 i.i.d.\ 가우시안(Gaussian)일 때, 프로베니우스 최소화는
% 최대 우도 추정(maximum likelihood estimation)과 동치이다.
% 그러나 큰 잔차에 민감하여 이상치(outlier)에 취약하다.

% \paragraph{$\ell_1$ 손실.}
% \begin{equation}
%   \mathcal{L}_1(\tensor{X}, \hat{\tensor{X}})
%   \;=\; \onenorm{\tensor{X} - \hat{\tensor{X}}}
%   \;=\; \sum_{i_1, \ldots, i_N}
%         \bigl|x_{i_1\cdots i_N} - \hat{x}_{i_1\cdots i_N}\bigr|.
%   \label{eq:loss_l1}
% \end{equation}
% $\mathcal{L}_1$ 최소화는 라플라스(Laplace) 잡음 모델에서의
% 최대 우도 추정에 해당하며, 양의 분해점(breakdown point)을 가지므로
% 희소 이상치에 강인(robust)하다.
% 그러나 원점에서 \emph{비매끄러운(non-smooth)} 함수로,
% 편미분이 $\partial|\cdot| = \sign(\cdot)$이므로
% 표준 최소제곱 풀이기를 직접 적용할 수 없다.

% \paragraph{Huber 손실.}
% Huber 손실~\cite{huber1964robust}은 원점 근방에서 $\ell_2$,
% 꼬리 부분에서 $\ell_1$처럼 작동하도록 매끄럽게 결합한 함수이다:
% \begin{equation}
%   \rho_\delta(r)
%   \;=\; \begin{cases}
%           \dfrac{1}{2} r^2
%             & |r| \leq \delta\text{인 경우,} \\[8pt]
%           \delta \!\left(|r| - \dfrac{\delta}{2}\right)
%             & |r| > \delta\text{인 경우,}
%         \end{cases}
%   \label{eq:huber}
% \end{equation}
% 여기서 $\delta > 0$은 전환 임계값이다.
% 집합 손실은
% $\mathcal{L}_H = \sum_{i_1,\ldots,i_N}
% \rho_\delta(x_{i_1\cdots i_N} - \hat{x}_{i_1\cdots i_N})$이다.
% $\mathcal{L}_H$는 모든 점에서 미분 가능하므로
% 부분기울기(subgradient) 없이 경사 하강법을 직접 적용할 수 있다.

% \paragraph{요약.}
% 표~\ref{tab:losses}는 세 손실 함수를 통계적 해석, 매끄러움 및
% 최적화 전략의 관점에서 비교한다.

% \begin{table}[h]
%   \centering
%   \caption{재구성 손실 함수 비교.}
%   \label{tab:losses}
%   \begin{tabular}{lccc}
%     \toprule
%     손실 함수 & 잡음 모델 & 매끄러움 & 최적화 방법 \\
%     \midrule
%     $\mathcal{L}_F$ (Frobenius)
%       & 가우시안       & $C^\infty$          & ALS (닫힌 형태) \\
%     $\mathcal{L}_1$ ($\ell_1$)
%       & 라플라스       & 원점에서 비매끄러움 & IRLS / 부분기울기 \\
%     $\mathcal{L}_H$ (Huber)
%       & 가우시안 + 이상치 & $C^1$            & 경사 하강법 \\
%     \bottomrule
%   \end{tabular}
% \end{table}
